\documentclass[11pt]{article}
\usepackage{amsmath, amssymb, amsthm, mathtools, bm}

\title{Prime-Indexed Recursive Tensor Mathematics with $\Lambda_m$-Governed Contraction:\\
Formal Specification, Examples, and Prior Art Positioning}
\author{Technical Defensive Publication}
\date{}

\begin{document}
\maketitle

\tableofcontents

\section{Executive Summary}

This document specifies a prime-indexed operator framework for discrete-time recursion on a tensor product space, governed by a scalar parameter $\Lambda_m$ which enforces contraction under explicit Lipschitz and operator norm constraints. The central result is a general contraction theorem for a class of $\Lambda_m$-interpolated update maps constructed from prime-indexed operator families and auxiliary linear components. The framework includes finite-dimensional diagonal and coupled linear examples, an infinite-dimensional generalization, and a basic implementation sketch in a numerical linear algebra environment.

The construction isolates the following elements:

\begin{itemize}
  \item A tensor product space $\mathcal{H}$ with a distinguished prime-indexed factor.
  \item A family of bounded linear operators $(A_p)_{p \in \mathbb{P}}$ on $\mathcal{H}$ indexed by prime numbers, aggregated into a bounded multiplicity operator $M$.
  \item An auxiliary operator $B$ that may depend linearly on both the state and an external input variable.
  \item A $\Lambda_m$-interpolated update map
  \[
    T_{\Lambda_m}(\Psi,x) = (1 - \Lambda_m)\Psi + \Lambda_m G(\Psi,x),
  \]
  where $G$ encodes the prime-indexed and auxiliary contributions.
  \item An explicit Lipschitz bound on $G$ in the state variable, yielding a sufficient condition for $T_{\Lambda_m}$ to be a contraction for suitable $\Lambda_m$.
\end{itemize}

Under the standard assumptions of completeness and boundedness, the contraction theorem implies existence and uniqueness of a fixed point for each frozen input and exponential convergence of iterates. The finite-dimensional examples demonstrate the mechanism in matrix form, including the case in which both the prime-indexed part and an additional linear term contribute to the Lipschitz constant.

This report is written entirely in formal mathematical terminology and is intended as a defensive publication specifying a precise operator class, its stability properties, and minimal implementation patterns.

\section{Prime-Indexed Recursive Tensor Space}

\subsection{Basic construction}

Let $\mathbb{P}$ denote the set of prime numbers and let $P \subset \mathbb{P}$ be either finite or countable. Define a separable Hilbert space $H_P$ with orthonormal basis $\{ e_p : p \in P \}$. Let $H_T$ be a separable Hilbert space (for instance $L^2(\mathbb{R})$ or $\ell^2(\mathbb{Z})$), and let $H_F = \mathbb{C}^d$ for some fixed $d \in \mathbb{N}$.

Define the tensor product space
\begin{equation}
  \mathcal{H} := H_P \widehat{\otimes} H_T \widehat{\otimes} H_F,
\end{equation}
equipped with the canonical Hilbert norm $\|\cdot\|$ induced by the tensor product. Let $(\mathcal{X}, \|\cdot\|_{\mathcal{X}})$ be a Banach space of input variables.

\subsection{Prime-indexed operators}

For each $p \in P$, let $A_p : \mathcal{H} \to \mathcal{H}$ be a bounded linear operator. Assume there is a sequence $(\alpha_p)_{p \in P} \subset [0,\infty)$ such that
\begin{equation}
  \|A_p\| \leq \alpha_p \quad \forall p \in P,
  \qquad
  \sum_{p \in P} \alpha_p < \infty.
\end{equation}
Then the series
\begin{equation}
  M := \sum_{p \in P} A_p
\end{equation}
converges in the operator norm and defines a bounded linear operator $M : \mathcal{H} \to \mathcal{H}$ with
\begin{equation}
  \|M\| \leq \sum_{p \in P} \alpha_p.
\end{equation}

\subsection{Auxiliary operator}

Let $B : \mathcal{H} \times \mathcal{X} \to \mathcal{H}$ be a map that is linear in its first argument and measurable in its second argument. In many cases of interest, $B$ will be of the form
\begin{equation}
  B(\Psi,x) = K \Psi + C x,
\end{equation}
where $K : \mathcal{H} \to \mathcal{H}$ and $C : \mathcal{X} \to \mathcal{H}$ are bounded linear operators.

\subsection{Prime-indexed recursive operator}

Define the prime-indexed recursive operator
\begin{equation}\label{eq:G_def}
  G(\Psi,x) := M\Psi + B(\Psi,x),
  \quad
  (\Psi,x) \in \mathcal{H} \times \mathcal{X}.
\end{equation}
This operator encodes the combined effect of the prime-indexed family $(A_p)_{p \in P}$ and the auxiliary term $B$.

\section{$\Lambda_m$-Governed Update and Stability Conditions}

\subsection{Definition of the $\Lambda_m$ update}

Let $\Lambda_m \in \mathbb{R}$ be a scalar parameter. Define
\begin{equation}\label{eq:T_Lm_def}
  T_{\Lambda_m}(\Psi,x)
    := (1 - \Lambda_m)\Psi
      + \Lambda_m G(\Psi,x),
\end{equation}
for $(\Psi,x) \in \mathcal{H} \times \mathcal{X}$. The corresponding recursion is
\begin{equation}
  \Psi_{t+1} = T_{\Lambda_m}(\Psi_t, x_t),
  \quad
  t = 0,1,2,\dots.
\end{equation}

\subsection{Lipschitz condition}

Assume that $G$ satisfies the following Lipschitz condition in its first argument.

\medskip

\noindent
\textbf{Assumption (Lipschitz bound).}
There exists a constant $L_G \in [0,1)$ such that for all $\Psi_1,\Psi_2 \in \mathcal{H}$ and all $x \in \mathcal{X}$,
\begin{equation}\label{eq:Lip_G}
  \|G(\Psi_1,x) - G(\Psi_2,x)\|
  \leq
  L_G \|\Psi_1 - \Psi_2\|.
\end{equation}
In particular, if $G$ has the form
\begin{equation}
  G(\Psi,x) = (M+K)\Psi + Cx
\end{equation}
with bounded $M,K,C$, then $L_G \leq \|M+K\|$.

\subsection{Effective contraction constant}

For any choice of $\Lambda_m \in \mathbb{R}$, define
\begin{equation}\label{eq:c_Lm_def}
  c(\Lambda_m) := (1 - \Lambda_m) + \Lambda_m L_G
  = 1 - \Lambda_m(1 - L_G).
\end{equation}
If $L_G < 1$ and $0 < \Lambda_m \leq 1$, then
\begin{equation}
  c(\Lambda_m) < 1.
\end{equation}

\section{General $\Lambda_m$ Contraction Theorem}

\subsection{Statement}

\begin{theorem}[$\Lambda_m$ contraction]
Let $(\mathcal{H},\|\cdot\|)$ be a Banach space and let $(\mathcal{X},\|\cdot\|_{\mathcal{X}})$ be a Banach space. Let $G : \mathcal{H} \times \mathcal{X} \to \mathcal{H}$ satisfy the Lipschitz bound \eqref{eq:Lip_G} for some $L_G \in [0,1)$. Let $\Lambda_m \in (0,1]$ and define $T_{\Lambda_m}$ by \eqref{eq:T_Lm_def}.

For any fixed $x \in \mathcal{X}$, define $T_{\Lambda_m}^x : \mathcal{H} \to \mathcal{H}$ by
\begin{equation}
  T_{\Lambda_m}^x(\Psi) := T_{\Lambda_m}(\Psi,x).
\end{equation}
Then:
\begin{enumerate}
  \item $T_{\Lambda_m}^x$ is a contraction with contraction constant $c(\Lambda_m)$ given by \eqref{eq:c_Lm_def}, i.e.,
  \[
    \|T_{\Lambda_m}^x(\Psi_1) - T_{\Lambda_m}^x(\Psi_2)\|
    \leq
    c(\Lambda_m) \|\Psi_1 - \Psi_2\|
    \quad \forall \Psi_1,\Psi_2 \in \mathcal{H}.
  \]
  \item There exists a unique $\Psi^\ast(x) \in \mathcal{H}$ such that
  \[
    \Psi^\ast(x) = T_{\Lambda_m}^x(\Psi^\ast(x)).
  \]
  \item For any initial $\Psi_0 \in \mathcal{H}$, the sequence $\{\Psi_t\}_{t \geq 0}$ defined by
  \[
    \Psi_{t+1} = T_{\Lambda_m}^x(\Psi_t)
  \]
  converges to $\Psi^\ast(x)$ and satisfies
  \[
    \|\Psi_t - \Psi^\ast(x)\|
    \leq
    c(\Lambda_m)^t \|\Psi_0 - \Psi^\ast(x)\|,
    \quad t = 0,1,2,\dots.
  \]
\end{enumerate}
\end{theorem}

\subsection{Proof}

For any $\Psi_1,\Psi_2 \in \mathcal{H}$ and fixed $x \in \mathcal{X}$,
\begin{align*}
  \|T_{\Lambda_m}^x(\Psi_1) - T_{\Lambda_m}^x(\Psi_2)\|
  &=
  \|(1 - \Lambda_m)(\Psi_1 - \Psi_2)
    + \Lambda_m (G(\Psi_1,x) - G(\Psi_2,x))\| \\
  &\leq
  (1 - \Lambda_m) \|\Psi_1 - \Psi_2\|
  + \Lambda_m \|G(\Psi_1,x) - G(\Psi_2,x)\| \\
  &\leq
  \bigl( (1 - \Lambda_m) + \Lambda_m L_G \bigr)
  \|\Psi_1 - \Psi_2\| \\
  &=
  c(\Lambda_m) \|\Psi_1 - \Psi_2\|.
\end{align*}
Since $L_G < 1$ and $\Lambda_m \in (0,1]$, it follows that $c(\Lambda_m) < 1$, so $T_{\Lambda_m}^x$ is a contraction on the complete metric space $(\mathcal{H},\|\cdot\|)$. The remaining conclusions follow from the Banach fixed point theorem.

\section{Finite-Dimensional Examples}

\subsection{Example 1: diagonal case}

Let $P = \{p_1,\dots,p_N\} \subset \mathbb{P}$. Let $H_P = \mathbb{C}^N$ with orthonormal basis $\{e_{p_k}\}_{k=1}^N$, and let $H_T = \mathbb{C}^m$, $H_F = \mathbb{C}^d$. Then
\begin{equation}
  \mathcal{H} := H_P \otimes H_T \otimes H_F
  \cong \mathbb{C}^{Nmd},
\end{equation}
with the Euclidean norm. Let $\mathcal{X} := \mathbb{C}^r$ with the Euclidean norm.

For each $p_k \in P$, choose $a_{p_k} \in \mathbb{C}$ and define $A_{p_k} : \mathcal{H} \to \mathcal{H}$ by
\begin{equation}
  A_{p_k}(e_{p_j} \otimes u \otimes v)
  :=
  \begin{cases}
    a_{p_k} e_{p_k} \otimes u \otimes v, & j = k, \\
    0, & j \neq k,
  \end{cases}
\end{equation}
for all $u \in H_T$, $v \in H_F$, extended linearly. Then $\|A_{p_k}\| = |a_{p_k}|$ and
\begin{equation}
  M := \sum_{k=1}^N A_{p_k}
\end{equation}
is a bounded diagonal operator with
\begin{equation}
  \|M\| \leq \sum_{k=1}^N |a_{p_k}|.
\end{equation}

Define $C : \mathcal{X} \to \mathcal{H}$ as a bounded linear map and set
\begin{equation}
  B(\Psi,x) := Cx.
\end{equation}
Then
\begin{equation}
  G(\Psi,x) := M\Psi + B(\Psi,x) = M\Psi + Cx
\end{equation}
satisfies
\begin{equation}
  \|G(\Psi_1,x) - G(\Psi_2,x)\|
  =
  \|M(\Psi_1 - \Psi_2)\|
  \leq
  \|M\| \|\Psi_1 - \Psi_2\|.
\end{equation}
Thus $L_G = \|M\|$. If $\|M\| < 1$ and $\Lambda_m \in (0,1]$, then $c(\Lambda_m) = (1 - \Lambda_m) + \Lambda_m \|M\| < 1$, and $T_{\Lambda_m}^x$ is a contraction for each fixed $x$.

\subsection{Example 2: coupled linear case}

Retain the same spaces. Let $M$ be as above and let $K : \mathcal{H} \to \mathcal{H}$ be a bounded linear operator. Define $C : \mathcal{X} \to \mathcal{H}$ bounded linear and set
\begin{equation}
  B(\Psi,x) := K\Psi + Cx.
\end{equation}
Then
\begin{equation}
  G(\Psi,x) := M\Psi + B(\Psi,x)
  =
  (M+K)\Psi + Cx.
\end{equation}
For any $\Psi_1,\Psi_2 \in \mathcal{H}$,
\begin{align*}
  \|G(\Psi_1,x) - G(\Psi_2,x)\|
  &=
  \|(M+K)(\Psi_1 - \Psi_2)\| \\
  &\leq
  \|M+K\| \|\Psi_1 - \Psi_2\|.
\end{align*}
Thus $L_G = \|M+K\|$. If $\|M+K\| < 1$ and $\Lambda_m \in (0,1]$, then
\begin{equation}
  c(\Lambda_m) = (1 - \Lambda_m) + \Lambda_m \|M+K\| < 1,
\end{equation}
and $T_{\Lambda_m}^x$ is a contraction on $\mathcal{H}$ for each fixed $x$.

\section{Infinite-Dimensional Generalization}

\begin{remark}[Infinite-dimensional extension]
Let $\mathcal{H}$ be a Banach or Hilbert space, let $M,K : \mathcal{H} \to \mathcal{H}$ be bounded linear operators, and let $C : \mathcal{X} \to \mathcal{H}$ be bounded linear. Define
\begin{equation}
  G(\Psi,x) := (M+K)\Psi + Cx,
  \quad
  T_{\Lambda_m}(\Psi,x)
  := (1 - \Lambda_m)\Psi + \Lambda_m G(\Psi,x).
\end{equation}
Then for fixed $x$,
\begin{equation}
  \|G(\Psi_1,x) - G(\Psi_2,x)\|
  \leq
  \|M+K\| \|\Psi_1 - \Psi_2\|,
\end{equation}
so $L_G = \|M+K\|$. The same contraction theorem applies if $\|M+K\| < 1$ and $\Lambda_m \in (0,1]$.

More generally, if direct control of $\|M+K\|$ is replaced by a spectral radius bound $\rho(M+K) < 1$ together with suitable norm inequalities (for instance via an equivalent norm in which the operator norm is bounded by a function of $\rho(M+K)$), then one obtains an effective $L_G < 1$ and recovers the contraction result with $c(\Lambda_m)$ as in \eqref{eq:c_Lm_def}.
\end{remark}

\section{Implementation Sketch}

\subsection{Finite-dimensional matrix representation}

In the finite-dimensional case $\mathcal{H} \cong \mathbb{C}^{n}$, one can represent $M$, $K$, and $C$ by matrices $M \in \mathbb{C}^{n \times n}$, $K \in \mathbb{C}^{n \times n}$, $C \in \mathbb{C}^{n \times r}$. The update
\begin{equation}
  T_{\Lambda_m}(\Psi,x) = (1 - \Lambda_m)\Psi + \Lambda_m\bigl((M+K)\Psi + Cx\bigr)
\end{equation}
reduces to a standard affine recursion.

\subsection{Code snippet}

A minimal numerical implementation can be written as follows.

\medskip

\noindent
\textbf{Pseudocode (Python-like).}
\begin{verbatim}
import numpy as np

def T_lambda_m(Psi, x, M, K, C, lambda_m):
    """
    Psi: state vector in C^n (as complex numpy array)
    x:   input vector in C^r
    M,K: n x n complex matrices
    C:   n x r complex matrix
    lambda_m: real scalar in (0,1]
    """
    G_Psi_x = (M + K) @ Psi + C @ x
    return (1.0 - lambda_m) * Psi + lambda_m * G_Psi_x

def iterate_T(Psi0, x, M, K, C, lambda_m, T):
    Psi = Psi0.copy()
    for _ in range(T):
        Psi = T_lambda_m(Psi, x, M, K, C, lambda_m)
    return Psi
\end{verbatim}

If the spectral norm of $M+K$ satisfies $\|M+K\|_2 < 1$ and $0 < \Lambda_m \leq 1$, then the matrix version of the contraction theorem applies with $L_G = \|M+K\|_2$ and $c(\Lambda_m)$ as in \eqref{eq:c_Lm_def}.

\section{Prior Art Positioning}

This framework specifies:

\begin{itemize}
  \item A tensor product state space with a distinguished prime-indexed factor.
  \item A prime-indexed family of bounded operators aggregated into a bounded multiplicity operator.
  \item A class of $\Lambda_m$-interpolated update maps composed with auxiliary linear components.
  \item A Lipschitz-based contraction theorem yielding fixed points and exponential convergence.
\end{itemize}

The combination of a prime-indexed operator family, explicit aggregation via an absolutely convergent series, and a $\Lambda_m$-governed convex interpolation with an abstract Lipschitz operator $G$ provides a concrete definition of a recursive operator system over a prime-indexed tensor space with explicit stability guarantees. The finite-dimensional diagonal and coupled examples illustrate the realization of such systems as matrix recursions with spectral norm bounds determining the contraction regime.

\appendix
\section*{Mathematical Appendix}
\addcontentsline{toc}{section}{Mathematical Appendix}

\section{Operator Norm Bounds for Prime-Indexed Aggregates}

\subsection{Boundedness of the multiplicity operator}

Let $P \subset \mathbb{P}$ be at most countable. For each $p \in P$ let
\[
  A_p : \mathcal{H} \to \mathcal{H}
\]
be a bounded linear operator. Assume the existence of a sequence
\[
  (\alpha_p)_{p \in P} \subset [0,\infty)
\]
such that
\begin{equation}\label{eq:A_p_bound_app}
  \|A_p\| \leq \alpha_p
  \quad \forall p \in P,
\end{equation}
and
\begin{equation}\label{eq:alpha_sum_app}
  \sum_{p \in P} \alpha_p < \infty.
\end{equation}

\begin{lemma}[Norm bound for $M = \sum_{p \in P} A_p$]
Under \eqref{eq:A_p_bound_app}--\eqref{eq:alpha_sum_app}, the series
\[
  M := \sum_{p \in P} A_p
\]
converges in the operator norm and $M : \mathcal{H} \to \mathcal{H}$ is bounded with
\begin{equation}
  \|M\| \leq \sum_{p \in P} \alpha_p.
\end{equation}
\end{lemma}

\begin{proof}
Let $(P_n)_{n \in \mathbb{N}}$ be an increasing sequence of finite subsets of $P$ with $\bigcup_n P_n = P$. Define
\[
  M_n := \sum_{p \in P_n} A_p.
\]
For $m > n$,
\[
  M_m - M_n = \sum_{p \in P_m \setminus P_n} A_p,
\]
so
\[
  \|M_m - M_n\|
  \leq
  \sum_{p \in P_m \setminus P_n} \|A_p\|
  \leq
  \sum_{p \in P_m \setminus P_n} \alpha_p.
\]
By \eqref{eq:alpha_sum_app}, the tail sums
\[
  \sum_{p \in P \setminus P_n} \alpha_p \to 0 \quad \text{as } n \to \infty.
\]
Hence $\{M_n\}$ is Cauchy in the operator norm and converges to some bounded operator $M$. For any $\Psi \in \mathcal{H}$,
\[
  \|M\Psi\|
  =
  \lim_{n \to \infty} \|M_n \Psi\|
  \leq
  \limsup_{n \to \infty}
  \left( \sum_{p \in P_n} \alpha_p \right)\|\Psi\|
  =
  \left( \sum_{p \in P} \alpha_p \right) \|\Psi\|.
\]
Thus $\|M\| \leq \sum_{p \in P} \alpha_p$.
\end{proof}

\subsection{Finite-dimensional diagonal case}

Consider the setting of Example~1 with
\[
  \mathcal{H} \cong \mathbb{C}^{Nmd}.
\]
For each $p_k \in P$ choose $a_{p_k} \in \mathbb{C}$ and define $A_{p_k}$ by
\[
  A_{p_k}(e_{p_j} \otimes u \otimes v)
  :=
  \begin{cases}
    a_{p_k} e_{p_k} \otimes u \otimes v, & j = k, \\
    0, & j \neq k.
  \end{cases}
\]
Then $A_{p_k}$ is diagonal in the $H_P$ component with a single non-zero block.

\begin{lemma}[Diagonal norm bound]
In the above finite-dimensional setting,
\[
  \|A_{p_k}\| = |a_{p_k}|,\quad
  \|M\| \leq \sum_{k=1}^N |a_{p_k}|.
\]
\end{lemma}

\begin{proof}
For any $u \in H_T$, $v \in H_F$,
\[
  \|A_{p_k}(e_{p_k} \otimes u \otimes v)\|
  =
  |a_{p_k}|\|u\|\|v\|.
\]
The operator $A_{p_k}$ acts as multiplication by $a_{p_k}$ on the subspace spanned by $e_{p_k} \otimes H_T \otimes H_F$ and vanishes on orthogonal subspaces. Thus $\|A_{p_k}\| = |a_{p_k}|$. The triangle inequality yields
\[
  \|M\| = \left\|\sum_{k=1}^N A_{p_k}\right\|
  \leq
  \sum_{k=1}^N \|A_{p_k}\|
  =
  \sum_{k=1}^N |a_{p_k}|.
\]
\end{proof}

\section{Lipschitz Bounds for $G(\Psi,x)$}

\subsection{Affine form in the state variable}

Let $M : \mathcal{H} \to \mathcal{H}$ and $K : \mathcal{H} \to \mathcal{H}$ be bounded linear operators, and let $C : \mathcal{X} \to \mathcal{H}$ be bounded linear. Define
\begin{equation}
  G(\Psi,x) := (M+K)\Psi + Cx.
\end{equation}

\begin{lemma}[Lipschitz constant for $G$]
For any $\Psi_1,\Psi_2 \in \mathcal{H}$ and $x \in \mathcal{X}$,
\[
  \|G(\Psi_1,x) - G(\Psi_2,x)\|
  \leq
  \|M+K\|\,\|\Psi_1 - \Psi_2\|.
\]
Hence $G$ is Lipschitz in the first argument with $L_G = \|M+K\|$.
\end{lemma}

\begin{proof}
We have
\[
  G(\Psi_1,x) - G(\Psi_2,x)
  =
  (M+K)(\Psi_1 - \Psi_2).
\]
Thus
\[
  \|G(\Psi_1,x) - G(\Psi_2,x)\|
  =
  \|(M+K)(\Psi_1 - \Psi_2)\|
  \leq
  \|M+K\|\,\|\Psi_1 - \Psi_2\|
\]
by the definition of the operator norm.
\end{proof}

\subsection{Decomposition into prime-indexed and auxiliary parts}

Suppose $M$ is constructed as the norm-convergent sum $M = \sum_{p \in P} A_p$ with $\|A_p\| \leq \alpha_p$ and $\sum\alpha_p < \infty$, and $K$ is any bounded linear operator. Then
\[
  \|M+K\| \leq \|M\| + \|K\| \leq \left( \sum_{p \in P} \alpha_p \right) + \|K\|.
\]
Thus an explicit Lipschitz bound is
\begin{equation}
  L_G \leq \sum_{p \in P} \alpha_p + \|K\|.
\end{equation}

\section{Proof of the General $\Lambda_m$ Contraction Theorem}

We give an explicit proof of the $\Lambda_m$ contraction theorem with the Lipschitz constant derived as above.

\subsection{Setup}

Let $(\mathcal{H},\|\cdot\|)$ be a Banach space and $(\mathcal{X},\|\cdot\|_{\mathcal{X}})$ a Banach space. Let $G : \mathcal{H} \times \mathcal{X} \to \mathcal{H}$ satisfy
\begin{equation}\label{eq:Lip_G_app}
  \|G(\Psi_1,x) - G(\Psi_2,x)\|
  \leq
  L_G \|\Psi_1 - \Psi_2\|
\end{equation}
for some $L_G \in [0,1)$ and for all $\Psi_1,\Psi_2 \in \mathcal{H}$, $x \in \mathcal{X}$. Fix $\Lambda_m \in (0,1]$ and define
\begin{equation}
  T_{\Lambda_m}(\Psi,x)
  :=
  (1 - \Lambda_m)\Psi + \Lambda_m G(\Psi,x).
\end{equation}
For any fixed $x \in \mathcal{X}$, define
\[
  T_{\Lambda_m}^x(\Psi) := T_{\Lambda_m}(\Psi,x).
\]

\subsection{Contraction constant}

\begin{lemma}
For any $\Psi_1,\Psi_2 \in \mathcal{H}$ and fixed $x \in \mathcal{X}$,
\begin{equation}
  \|T_{\Lambda_m}^x(\Psi_1) - T_{\Lambda_m}^x(\Psi_2)\|
  \leq
  c(\Lambda_m) \|\Psi_1 - \Psi_2\|,
\end{equation}
where
\begin{equation}
  c(\Lambda_m) := (1 - \Lambda_m) + \Lambda_m L_G
  = 1 - \Lambda_m(1 - L_G).
\end{equation}
If $L_G < 1$ and $0 < \Lambda_m \leq 1$, then $c(\Lambda_m) < 1$.
\end{lemma}

\begin{proof}
We compute:
\begin{align*}
  T_{\Lambda_m}^x(\Psi_1) - T_{\Lambda_m}^x(\Psi_2)
  &=
  (1 - \Lambda_m)(\Psi_1 - \Psi_2)
  + \Lambda_m (G(\Psi_1,x) - G(\Psi_2,x)).
\end{align*}
Hence
\begin{align*}
  \|T_{\Lambda_m}^x(\Psi_1) - T_{\Lambda_m}^x(\Psi_2)\|
  &\leq
  (1 - \Lambda_m) \|\Psi_1 - \Psi_2\|
  + \Lambda_m \|G(\Psi_1,x) - G(\Psi_2,x)\| \\
  &\leq
  (1 - \Lambda_m) \|\Psi_1 - \Psi_2\|
  + \Lambda_m L_G \|\Psi_1 - \Psi_2\| \\
  &=
  \left( (1 - \Lambda_m) + \Lambda_m L_G \right)
  \|\Psi_1 - \Psi_2\| \\
  &=
  c(\Lambda_m) \|\Psi_1 - \Psi_2\|.
\end{align*}
If $L_G < 1$ and $0 < \Lambda_m \leq 1$, then $1 - L_G > 0$ and
\[
  c(\Lambda_m) = 1 - \Lambda_m(1 - L_G) < 1.
\]
\end{proof}

\subsection{Fixed point existence and uniqueness}

\begin{lemma}
For each fixed $x \in \mathcal{X}$, the map $T_{\Lambda_m}^x$ has a unique fixed point $\Psi^\ast(x) \in \mathcal{H}$.
\end{lemma}

\begin{proof}
Since $T_{\Lambda_m}^x$ is a contraction on the complete metric space $(\mathcal{H},\|\cdot\|)$ with contraction constant $c(\Lambda_m) < 1$, the Banach fixed point theorem implies the existence of a unique $\Psi^\ast(x) \in \mathcal{H}$ such that
\[
  \Psi^\ast(x) = T_{\Lambda_m}^x(\Psi^\ast(x)).
\]
\end{proof}

\subsection{Exponential convergence}

\begin{lemma}
Let $x \in \mathcal{X}$ be fixed and let $\Psi^\ast(x)$ be the fixed point of $T_{\Lambda_m}^x$. For any initial $\Psi_0 \in \mathcal{H}$, the iterates
\[
  \Psi_{t+1} = T_{\Lambda_m}^x(\Psi_t)
\]
satisfy
\[
  \|\Psi_t - \Psi^\ast(x)\|
  \leq
  c(\Lambda_m)^t \|\Psi_0 - \Psi^\ast(x)\|
  \quad \text{for all } t \in \mathbb{N}_0.
\]
\end{lemma}

\begin{proof}
By induction:
\[
  \|\Psi_{t+1} - \Psi^\ast(x)\|
  =
  \|T_{\Lambda_m}^x(\Psi_t) - T_{\Lambda_m}^x(\Psi^\ast(x))\|
  \leq
  c(\Lambda_m)\|\Psi_t - \Psi^\ast(x)\|.
\]
Since $\Psi_0$ is given, repeated application yields
\[
  \|\Psi_t - \Psi^\ast(x)\|
  \leq
  c(\Lambda_m)^t \|\Psi_0 - \Psi^\ast(x)\|.
\]
\end{proof}

\section{Spectral Radius Considerations}

\subsection{Norm versus spectral radius}

Let $L : \mathcal{H} \to \mathcal{H}$ be a bounded linear operator. Its spectral radius is
\[
  \rho(L) := \sup\{ |\lambda| : \lambda \in \sigma(L)\}.
\]
In general, $\rho(L) \leq \|L\|$. In finite dimensions, one may use equivalence of norms to obtain an equivalent norm $\|\cdot\|'$ such that
\[
  \|L\|' < 1
\quad \text{if} \quad
  \rho(L) < 1.
\]

\begin{lemma}
Let $L$ be a bounded linear operator on a finite-dimensional vector space. If $\rho(L) < 1$, then there exists a norm $\|\cdot\|'$ such that $\|L\|' < 1$.
\end{lemma}

\begin{proof}
In finite dimensions the Jordan normal form or the Schur decomposition may be used to construct a basis in which $L$ is upper triangular with diagonal entries strictly inside the unit disk. One can then define a norm adapted to this basis in which the operator norm of $L$ is strictly less than $1$. Details are standard and omitted.
\end{proof}

\subsection{Effective Lipschitz constant from spectral control}

In contexts where $G(\Psi,x) = L\Psi + Cx$, with bounded $L$ and $C$, and one knows a bound $\rho(L) < 1$ but not a tight norm bound, one may argue as follows.

\begin{lemma}
Let $L$ be a bounded linear operator on a finite-dimensional space with $\rho(L) < 1$. Then there exists a norm $\|\cdot\|'$ and a constant $L_G' < 1$ such that
\[
  \|L\Psi\|' \leq L_G' \|\Psi\|'
  \quad \forall \Psi.
\]
In this norm, $G(\Psi,x) = L\Psi + Cx$ is Lipschitz in $\Psi$ with Lipschitz constant $L_G'$, and the $\Lambda_m$ contraction theorem holds with $L_G'$ in place of $L_G$ and $c(\Lambda_m) = (1 - \Lambda_m) + \Lambda_m L_G'$.
\end{lemma}

\begin{proof}
Apply the previous lemma to obtain a norm $\|\cdot\|'$ with $\|L\|' < 1$, and set $L_G' := \|L\|'$.
\end{proof}

This yields a spectral-radius-based route to contraction, complementing the direct norm-based conditions.

\section{Summary of Norm and Contraction Conditions}

For convenience, we summarize the explicit conditions under which the $\Lambda_m$-governed update
\[
  T_{\Lambda_m}(\Psi,x)
  =
  (1 - \Lambda_m)\Psi + \Lambda_m G(\Psi,x)
\]
is a contraction in the state variable for fixed input.

\begin{itemize}
  \item If $G(\Psi,x) = (M+K)\Psi + Cx$ with bounded $M,K,C$ and
    \[
      \|M+K\| < 1,
      \quad \Lambda_m \in (0,1],
    \]
    then $G$ is Lipschitz with $L_G = \|M+K\|$ and
    \[
      c(\Lambda_m) = (1 - \Lambda_m) + \Lambda_m \|M+K\| < 1.
    \]
  \item More generally, if $G$ satisfies
    \[
      \|G(\Psi_1,x) - G(\Psi_2,x)\|
      \leq L_G \|\Psi_1 - \Psi_2\|
      \quad \text{for some } L_G < 1,
    \]
    then for any $\Lambda_m \in (0,1]$,
    \[
      c(\Lambda_m) = 1 - \Lambda_m(1 - L_G) < 1,
    \]
    and $T_{\Lambda_m}^x$ is a contraction for each $x$.
  \item In finite dimensions, if $G(\Psi,x) = L\Psi + Cx$ and $\rho(L) < 1$, then an equivalent norm exists in which $\|L\|' < 1$, yielding a Lipschitz constant $L_G' < 1$ and contraction in that norm.
\end{itemize}

These bounds complete the explicit mathematical characterization of the $\Lambda_m$-governed recursion and its stability properties.


\nocite{*}
\bibliographystyle{plain}
\bibliography{references}
\end{document}